% !TeX root = ../thesis.tex
\chapter{知识抽取模型工程化部署方法}

\section{工程化部署总体架构}
模型工程化部署的核心目标，是将上一章形成的模型制品稳定接入部署平台，并形成可发布、可回归、可回滚的服务链路。为此，本文采用“模型制品层、推理服务层、部署管理层、监控治理层”四层架构。模型制品层负责统一管理 ONNX 文件和标签元数据；推理服务层负责提供统一输入输出接口；部署管理层负责在 GPUStack 中完成实例启动、Profile 绑定和节点调度；监控治理层负责输出日志、指标与回滚依据。该架构在实现上采用 GPUStack 的部署组织方式，并结合 ONNX Runtime 的推理运行能力组织统一服务链路。

\begin{figure}[htbp]
  \centering
  \begin{tikzpicture}[
    box/.style={draw, rectangle, text width=3.6cm, minimum height=0.95cm, align=center, inner sep=3pt, font=\wuhao},
    arr/.style={-Latex, semithick}
  ]
    \node[box] (a) at (0,0) {模型制品层};
    \node[box] (b) at (0,-1.45) {统一推理服务层};
    \node[box] (c) at (0,-2.90) {GPUStack 部署层};
    \node[box] (d) at (0,-4.35) {监控治理层};
    \draw[arr] (a) -- (b);
    \draw[arr] (b) -- (c);
    \draw[arr] (c) -- (d);
  \end{tikzpicture}
  \caption{模型工程化部署总体架构}
  \label{fig:deploy_arch_compact}
\end{figure}

\section{统一推理服务实现}
统一推理服务实现的目标，是在不改变平台调用方式的前提下，对不同知识抽取模型提供一致的服务入口、返回结构与版本标识，从而使模型切换、实例编排和后续治理均可在统一接口之上完成。围绕这一目标，本文重点设计请求响应协议、执行流程以及与部署侧的衔接方式，并将其作为后续本地化部署和性能优化的共同基础。

为了避免平台侧为每个模型单独适配接口，本文在服务层定义统一请求格式。设请求对象为
\begin{equation}
q=\langle x, s, m, p \rangle,
\label{eq:request_compact}
\end{equation}
其中 $x$ 为输入文本，$s$ 为场景标识，$m$ 为模型编号，$p$ 为部署 Profile。服务输出定义为
\begin{equation}
y=\langle \hat{\mathcal{T}}, c, \ell, v \rangle,
\label{eq:response_compact}
\end{equation}
其中 $\hat{\mathcal{T}}$ 为三元组结果，$c$ 为整体置信度，$\ell$ 为推理时延，$v$ 为模型版本。该设计保证了模型替换不会影响平台侧调用协议。

\begin{table}[htbp]
  \centering
  \caption{统一推理服务的关键接口字段}
  \vspace{0.5em}\wuhao
  \begin{tabularx}{\textwidth}{l X}
    \toprule
    接口字段 & 字段说明 \\
    \midrule
    \texttt{text} & 待抽取文本内容 \\
    \texttt{scene} & 场景标识，如工艺文件、故障报告、标准规范 \\
    \texttt{model\_id} & 指定调用模型，支持显式选择和系统推荐 \\
    \texttt{profile} & 指定部署配置，用于区分延迟优先、吞吐优先等策略 \\
    \texttt{triples} & 统一结构化输出结果，包含实体、关系和证据片段 \\
    \texttt{version} & 当前模型版本，用于追踪与回滚 \\
    \bottomrule
  \end{tabularx}
\label{tab:service_schema_compact}
\end{table}

\section{本地化部署实现}
在本地化部署实现中，本文保留 GPUStack 的基础能力，包括模型文件上传、模型实例创建、节点管理、基准测试与结果导出，而不引入大模型推理、远程模型市场等与本文目标无关的复杂功能。这样做的主要原因在于，知识抽取模型库的部署对象是统一导出的 ONNX 制品，其核心需求是稳定部署、性能对比和系统演示，而不是开放式生成服务。

从部署流程看，系统首先接收本地上传的模型制品包，随后解析元数据、标签配置与默认 Profile，并在部署页面中完成模型、节点、精度模式和推理后端的选择。实例启动后，系统通过健康检查确认接口状态，再将实例状态、运行日志和资源指标反馈到前端页面。上述过程保证了模型部署不再依赖人工命令行操作，而能够被统一纳入系统管理流程。

\section{部署优化方法}
针对导师提出的“编写几个 Profile 文件进行模型部署优化”要求，本文不将 Profile 简单理解为零散部署参数，而是将其定义为配置驱动的性能 profiling 文件与部署模板。具体而言，Profile 定义为
\begin{equation}
\mathbf{p}=[b,c,l,e,r],
\label{eq:profile_compact}
\end{equation}
其中 $b$ 表示 batch size，$c$ 表示并发数，$l$ 表示最大序列长度，$e$ 表示精度模式，$r$ 表示推理后端。任一 Profile 的可执行性必须满足显存约束：
\begin{equation}
M_{\text{static}}+M_{\text{act}}(\mathbf{p})+M_{\text{ws}}(\mathbf{p}) \leq \rho M_{\text{gpu}},
\label{eq:profile_memory_compact}
\end{equation}
其中 $\rho$ 为安全冗余系数。为兼顾质量与性能，本文采用如下综合评分函数：
\begin{equation}
\mathcal{S}(m,\mathbf{p})=\lambda_1 F_1(m)+\lambda_2 T(m,\mathbf{p})-\lambda_3 L_{95}(m,\mathbf{p})-\lambda_4 M(m,\mathbf{p}),
\label{eq:profile_score_compact}
\end{equation}
其中 $T$ 为吞吐量，$L_{95}$ 为 P95 延迟，$M$ 为显存开销。

\begin{table}[htbp]
  \centering
  \caption{面向知识抽取模型的典型 Profile 模板}
  \vspace{0.5em}\wuhao
  \begin{tabularx}{\textwidth}{l c c c c X}
    \toprule
    模板名称 & Batch & 并发数 & 最大长度 & 精度模式 & 适用场景 \\
    \midrule
    latency-L1 & 1 & 1 & 128 & FP16 & 实时接口、低时延优先 \\
    throughput-T1 & 4 & 8 & 128 & FP16 & 批量补录、吞吐量优先 \\
    memory-M1 & 1 & 2 & 128 & FP16 & 资源受限环境 \\
    regression-R1 & 2 & 4 & 256 & FP16 & 发布前回归验证 \\
    \bottomrule
  \end{tabularx}
  \label{tab:profile_templates_compact}
\end{table}

在原型系统的工程材料中，上述模板进一步落地为 \texttt{profiles/default.yaml}、\texttt{profiles/latency-l1.yaml}、\texttt{profiles/throughput-t1.yaml}、\texttt{profiles/memory-m1.yaml} 和 \texttt{profiles/regression-r1.yaml} 五类配置文件，分别对应默认部署、低延迟服务、高吞吐批处理、资源受限部署与版本回归验证场景。这样一来，Profile 不再只是论文中的抽象概念，而成为可直接驱动 benchmark 任务与部署比较的配置对象。

\subsection{优化目标与指标体系}
面向知识抽取模型的性能优化，并不是简单追求某一项指标的最大化。对于在线故障报告抽取场景，优先目标通常是较低延迟和稳定返回；对于离线工艺知识补录场景，优先目标通常是较高吞吐量；对于企业资源紧张环境，则更强调显存可控和实例部署成功率。因此，本文将部署优化目标划分为延迟优先、吞吐优先和资源优先三类，并通过统一指标体系进行比较。

\subsection{Profile 模板与选择方法}
\begin{algorithm}[htbp]
\caption{面向部署场景的 Profile 选择方法}
\KwIn{模型集合 $\mathcal{M}$，Profile 集合 $\mathcal{P}$，场景约束 $\mathbf{s}$}
\KwOut{推荐模型与推荐 Profile}
初始化最优得分 $score^{*}\leftarrow -\infty$\;
\ForEach{$m\in\mathcal{M}$}{
  \ForEach{$\mathbf{p}\in\mathcal{P}$}{
    判断 $\mathbf{p}$ 是否满足式~(\ref{eq:profile_memory_compact})\;
    \If{可执行}{
      计算综合得分 $\mathcal{S}(m,\mathbf{p})$\;
      \If{$\mathcal{S}(m,\mathbf{p})>score^{*}$}{
        更新推荐模型与 Profile\;
      }
    }
  }
}
\end{algorithm}

\section{发布与运维治理}
\subsection{发布回归与版本门槛}
工程化部署的关键，不只是“把模型部署起来”，还包括“如何稳定发布”和“如何快速回滚”。本文为每个模型版本保留稳定基线，并将发布门槛定义为
\begin{equation}
\Delta F_{1}\geq -\epsilon_{1},\quad \Delta L_{95}\leq \epsilon_{2},\quad \Delta M\leq \epsilon_{3},
\label{eq:release_gate}
\end{equation}
即新版本在精度、延迟和显存上不能出现超出阈值的退化。若门槛不满足，则系统自动回退至稳定版本。

\begin{figure}[htbp]
  \centering
  \begin{tikzpicture}[
    box/.style={draw, rectangle, text width=2.35cm, minimum height=0.95cm, align=center, inner sep=3pt, font=\wuhao},
    arr/.style={-Latex, semithick}
  ]
    \node[box] (a) at (0,0) {模型登记};
    \node[box] (b) at (3.0,0) {绑定 Profile};
    \node[box] (c) at (6.0,0) {部署测试实例};
    \node[box] (d) at (9.0,0) {健康检查与压测};
    \node[box] (e) at (12.0,0) {发布或回滚};
    \draw[arr] (a) -- (b);
    \draw[arr] (b) -- (c);
    \draw[arr] (c) -- (d);
    \draw[arr] (d) -- (e);
  \end{tikzpicture}
  \caption{模型版本发布与回滚流程}
  \label{fig:deploy_release_compact}
\end{figure}

\begin{table}[htbp]
  \centering
  \caption{模型版本发布检查项}
  \vspace{0.5em}\wuhao
  \begin{tabularx}{\textwidth}{l X}
    \toprule
    检查项 & 检查内容 \\
    \midrule
    制品完整性 & ONNX 文件、标签映射、元数据、默认 Profile 是否齐全 \\
    接口稳定性 & 服务输入输出字段是否符合统一 schema \\
    性能回归 & 新版本在延迟、吞吐量、显存上的变化是否在阈值内 \\
    业务样例核查 & 关键工艺、故障和标准文本样例输出是否保持一致 \\
    日志可追踪性 & 能否定位到具体模型版本、节点和 Profile \\
    \bottomrule
  \end{tabularx}
\label{tab:release_gate_compact}
\end{table}

\subsection{可观测与回滚机制}

\begin{figure}[htbp]
  \centering
  \begin{tikzpicture}[
    box/.style={draw, rectangle, text width=2.55cm, minimum height=0.95cm, align=center, inner sep=3pt, font=\wuhao},
    arr/.style={-Latex, semithick}
  ]
    \node[box] (a) at (0,1.1) {推理日志};
    \node[box] (b) at (0,-1.1) {性能指标};
    \node[box] (c) at (4.2,0) {告警规则};
    \node[box] (d) at (8.4,0) {回滚决策};
    \draw[arr] (a) -- (c);
    \draw[arr] (b) -- (c);
    \draw[arr] (c) -- (d);
  \end{tikzpicture}
  \caption{部署可观测与回滚联动机制}
  \label{fig:observability_loop_compact}
\end{figure}

\section{本章小结}
本章完成了模型工程化部署方案设计。通过统一推理接口、Profile 驱动优化、发布回归和回滚机制，本文将知识抽取模型库从“可训练模型集合”推进为“可治理模型服务集合”，为下一章的系统实现与应用验证提供了可运行基础。
